\begin{enumerate}
    \item El banco presenta una tasa de mora del
12.
    \item 7%
, equivalente a 127 clientes de una cartera de 1,000.
    \item Este nivel de incumplimiento representa un indicador que debe monitorearse de forma continua, especialmente ante un posible crecimiento de la cartera de clientes.
    \item La distribución de los clientes con mora entre los distintos sectores económicos es relativamente homogénea, por lo que el sector laboral no evidencia un comportamiento diferenciador que permita explicar el incumplimiento de pago.
    \item La relación entre el nivel de ingreso y el monto del préstamo no presenta un patrón lineal claramente definido, lo que indica que estas variables, por sí solas, no permiten identificar diferencias evidentes entre clientes con y sin mora.
    \item La antigüedad laboral muestra una distribución similar para ambos grupos de clientes, por lo que tampoco constituye un factor con suficiente capacidad explicativa cuando se analiza de manera individual.
    \item La mayor concentración de clientes en mora se encuentra en los rangos de edad comprendidos entre los
20 y 64 años.
    \item Sin embargo, el análisis descriptivo no permite concluir que la edad sea un factor determinante del riesgo crediticio, ya que la distribución observada también puede estar influenciada por la composición de la cartera de clientes.
\end{enumerate}
